%! The Language of Variation — Unit 1, Lesson 1.1 — Guided Notes & Practice (Answer Key)
% MIRROR of ../notes/main.tex: -key in place of -boxes, \termblank -> \termans, \blank -> \ans, and the
% answer argument of every \pcell, \writespace and \labelbox filled. Nothing else differs.
% THE in-class document, in the MAIN IDEAS / NOTES shape (LESSON_SHAPE.md §1): the vocabulary
% box, then ONE two-column guidednotes table. Each row is one idea — a label on the left; on the
% right a SHORT printed definition, ONE large pre-drawn display, then a grid of a few numbered
% problems with real writing room. Problems are numbered 1..N through the component.
% DENSITY: a \blank{} only where a word or number IS the answer (the C/Q strip in the Amounts
% row); no mid-sentence blanks; two problems across; 2–3cm of answer space each.
%   I DO   : the teacher reads the definition, marks up the display, works the first problem.
%   WE DO  : the class works the rest of each grid, students holding the pen; the last row,
%            Guided Practice, is worked entirely together in a fresh context.
%   YOU DO : nothing on this page — ../homework/ IS the individual practice.
% The crux is the AMOUNTS row — a column written entirely in digits is not automatically
% quantitative. Permit numbers, parking spots, and ZIP codes average cleanly and mean nothing
% (EK VAR-1.C.1, VAR-1.C.2).
%
% THE CONTEXTS: three hiking parties (Trail / Time out), two permits with their water bottles,
% the ranger's Short / Medium / Long regrouping, the five-column digits strip, and the Season
% Pass for Guided Practice.
%
% PAGE LOCKSTEP with notes_key/: \termblank -> \termans, \blank -> \ans, and the answer argument
% of every \pcell, \writespace and \labelbox filled. Nothing else differs. A row cannot break
% across pages: figure and grid are separate sub-rows (`\\` then `& ...`).
\documentclass[12pt]{article}
\usepackage{apstats-article}
\usepackage{apstats-key}   % pulls in -boxes; do NOT also load -boxes

\begin{document}

\pageheader{Unit 1, Lesson 1.1}{Guided Notes \& Practice --- Answer Key}
% NAMESTRIP: no \namedateperiod here — the name row lives on the cover only.

\vspace{2pt}

\boxguard[16]
\begin{vocabbox}
\small
Fill in each term \textbf{as we name it} in the notes below.
\par
\termans{Categorical variable}{takes values that are category names or group labels --- Gold, Cedar Falls, Yes.}
\par
\termans{Quantitative variable}{takes numerical values for a measured or counted amount --- minutes, pounds, bottles.}
\par
\termans{The amount test}{ask whether the value records how much or how many, not whether it is written in digits.}
\par
\termans{Grouping (binning)}{replacing measured amounts with group labels; turns Q into C, and the mean cannot be recovered.}
\par
\end{vocabbox}

\small
\begin{guidednotes}
% ── ROW 1: two kinds of variable ─────────────────────────────────────────────
\mainidea[Two kinds of]{Variable} &
A \textbf{variable} is anything recorded for each individual. \textbf{Categorical} values are
\emph{group labels}; \textbf{quantitative} values are \emph{numbers for a measured or counted
amount}.

\vspace{4pt}
\begin{center}
\begin{tikzpicture}[scale=1.0]
  % column shading
  \fill[goldbg] (2.3,-0.3) rectangle (5.6,2.9);
  \fill[greenbg] (5.6,-0.3) rectangle (9.2,2.9);
  % header
  \fill[navy] (0,2.2) rectangle (9.2,2.9);
  \node[white, font=\bfseries\small] at (1.15,2.55) {Party};
  \node[white, font=\bfseries\small] at (3.95,2.55) {Trail};
  \node[white, font=\bfseries\small] at (7.4,2.55) {Time out (min)};
  % rows
  \foreach \y/\p/\t/\m in {1.6/Alvarez/Ridge Loop/95, 0.9/Boone/Cedar Falls/140,
                           0.2/Delgado/Summit Trail/215} {
    \node[font=\small] at (1.15,\y) {\p};
    \node[font=\small] at (3.95,\y) {\t};
    \node[font=\small] at (7.4,\y) {\m};}
  \draw[linegray] (0,-0.3) rectangle (9.2,2.9);
  \draw[linegray] (2.3,-0.3) -- (2.3,2.9);
  \draw[linegray] (5.6,-0.3) -- (5.6,2.9);
  % arrows to the label boxes
  \draw[->, thick, goldacc] (3.95,-0.35) -- (3.95,-0.95);
  \draw[->, thick, greenacc] (7.4,-0.35) -- (7.4,-0.95);
  \node[anchor=north] at (3.95,-1.0) {\labelbox{2.6cm}{categorical}};
  \node[anchor=north] at (7.4,-1.0) {\labelbox{2.6cm}{quantitative}};
  \node[font=\scriptsize\itshape, text width=3.2cm, align=center, anchor=north] at (3.95,-2.05)
    {Summit Trail is not ``more'' than Ridge Loop};
  \node[font=\scriptsize\itshape, text width=3.2cm, align=center, anchor=north] at (7.4,-2.05)
    {215 really is 120 more than 95};
\end{tikzpicture}
\end{center}
\\
& \notesprompt{Categorical or quantitative? Say why.}
\begin{probgrid}{|Y|Y|}
\pcell{1}{Group size (number of hikers)}{1.8cm}{Q --- a count is an amount; 6 hikers is more than 2} &
\pcell{2}{Dog along (Yes / No)}{1.8cm}{C --- Yes and No are labels} \\ \hline
\end{probgrid}
\\
& \begin{probgrid}*{|Y|Y|}
\pcell{3}{Pack weight (pounds)}{1.8cm}{Q --- a measured amount of weight} &
\pcell{4}{Weather at the trailhead}{1.8cm}{C --- Sunny, Rain, Fog: group labels} \\ \hline
\end{probgrid}
\\ \hline
% ── ROW 2: the amount test ───────────────────────────────────────────────────
\mainidea[The]{Amount test} &
Do not ask whether it has digits. Ask: \textbf{is the value an \emph{amount} of something ---
how much, or how many?} Second form: \textbf{does the average mean anything?}

\vspace{4pt}
\begin{center}
\begin{tikzpicture}[scale=1.0]
  % ── permits ──
  \node[font=\bfseries\small\color{navy}] at (1.9,3.1) {Permit \#};
  \foreach \x/\n in {0.9/2041, 2.9/6318} {
    \fill[sky, rounded corners=2pt] (\x-0.75,1.8) rectangle (\x+0.75,2.75);
    \draw[navy, rounded corners=2pt] (\x-0.75,1.8) rectangle (\x+0.75,2.75);
    \node[font=\tiny\color{slate}] at (\x,2.52) {PERMIT};
    \node[font=\bfseries\normalsize] at (\x,2.12) {\n};}
  \draw[->, thick] (1.9,1.7) -- (1.9,1.15);
  \node[font=\small] at (1.9,0.85) {average $=4179.5$};
  \node[font=\bfseries\large\color{redacc}] at (1.9,0.25) {nobody's permit};
  \node[font=\small] at (1.6,-0.45) {Permit \# is \labelbox{2.2cm}{categorical}};
  % ── bottles ──
  \node[font=\bfseries\small\color{navy}] at (6.6,3.1) {Water bottles};
  \foreach \xo/\n in {4.6/6, 7.3/2} {
    \foreach \k in {1,...,\n} {
      \fill[navy!55, rounded corners=1pt] ({\xo+0.32*\k-0.1},1.85) rectangle ({\xo+0.32*\k+0.1},2.35);
      \fill[navy] ({\xo+0.32*\k-0.05},2.35) rectangle ({\xo+0.32*\k+0.05},2.5);}}
  \node[font=\small] at (5.75,2.7) {6};
  \node[font=\small] at (7.8,2.7) {2};
  \draw[->, thick] (6.6,1.7) -- (6.6,1.15);
  \node[font=\small] at (6.6,0.85) {average $=4$};
  \foreach \k in {1,...,4} {
    \fill[navy!55, rounded corners=1pt] ({5.9+0.32*\k-0.1},0.0) rectangle ({5.9+0.32*\k+0.1},0.5);
    \fill[navy] ({5.9+0.32*\k-0.05},0.5) rectangle ({5.9+0.32*\k+0.05},0.65);}
  \node[font=\small] at (7.2,-0.45) {Water bottles is \labelbox{2.2cm}{quantitative}};
\end{tikzpicture}
\end{center}

\vspace{2pt}
\textbf{In your own words:} the amount test.
\writespace{1.6cm}{Is the value an amount --- how much or how many? If yes, Q; if it only names a group, C.}
\\
& \notesprompt{Run the test.}
\begin{probgrid}{|Y|Y|}
\pcell{5}{Four water bottles --- what real thing is that?}{2.0cm}{A real amount: between them the parties carry 4 bottles each. Quantitative.} &
\pcell{6}{Hand me the party whose permit is 4179.5.}{2.0cm}{Nobody --- no such party. A permit number only names a party. Categorical.} \\ \hline
\end{probgrid}
\\ \hline
% ── ROW 3: grouping runs one way ─────────────────────────────────────────────
\mainidea[Grouping is]{One way only} &
\textbf{Grouping (binning)} replaces amounts with group labels. It runs one way: 95 minutes
becomes \emph{Short}, but \emph{Short} never becomes 95 again.

\vspace{4pt}
\begin{center}
\begin{tikzpicture}[scale=1.0]
  % ruler
  \fill[sky] (0,1.5) rectangle (4.0,2.1);
  \fill[goldbg] (4.0,1.5) rectangle (6.0,2.1);
  \fill[greenbg] (6.0,1.5) rectangle (8.0,2.1);
  \draw[thick] (0,1.5) rectangle (8.0,2.1);
  \draw[thick] (4.0,1.5) -- (4.0,2.1); \draw[thick] (6.0,1.5) -- (6.0,2.1);
  \foreach \x/\l in {0/0, 2/60, 4/120, 6/180, 8/240} {
    \draw (\x,1.5) -- (\x,1.35); \node[below, font=\tiny] at (\x,1.35) {\l};}
  \node[font=\tiny] at (4.0,0.75) {\emph{Time out}, minutes};
  % the three hikes
  \foreach \x/\p in {3.167/95, 4.667/140, 7.167/215} {
    \fill[navy] (\x,1.8) circle (0.11);
    \node[above, font=\scriptsize\bfseries] at (\x,2.15) {\p};}
  % one-way arrow down to the label column
  \draw[->, very thick] (4.3,0.55) -- (4.3,-0.2);
  \node[right, font=\scriptsize\itshape] at (4.45,0.18) {amounts $\rightarrow$ labels};
  \draw[->, very thick, redacc] (3.5,-0.2) -- (3.5,0.55);
  \draw[redacc, very thick] (3.3,0.0) -- (3.7,0.4);
  \draw[redacc, very thick] (3.3,0.4) -- (3.7,0.0);
  \node[left, font=\scriptsize\itshape\color{redacc}] at (3.35,0.18) {never back};
  \foreach \x/\l/\c in {2.0/Short/sky, 5.0/Medium/goldbg, 7.0/Long/greenbg} {
    \fill[\c] (\x-0.9,-0.95) rectangle (\x+0.9,-0.3);
    \draw[thick] (\x-0.9,-0.95) rectangle (\x+0.9,-0.3);
    \node[font=\bfseries\small] at (\x,-0.62) {\l};}
  \node[font=\tiny] at (4.0,-1.25) {the new \emph{Length} column};
  \node[font=\small] at (4.0,-1.85) {The new column is \labelbox{2.6cm}{categorical}};
\end{tikzpicture}
\end{center}
\\
& \notesprompt{Use the two columns.}
\begin{probgrid}{|Y|Y|}
\pcell{7}{Which column can be averaged? Compute it.}{2.0cm}{Time out: $(95+140+215)/3 = 150$ minutes} &
\pcell{8}{Boone now reads \emph{Medium}. Can the ranger get 140 back? Why?}{2.0cm}{No. Medium covers every time from 120 to 180; the exact minutes are gone.} \\ \hline
\end{probgrid}
\\
& \begin{probgrid}*{|Y|Y|}
\pcell{9}{Name one thing the ranger can no longer compute.}{2.0cm}{The mean time out (or the longest hike, or the total minutes)} &
\pcell{10}{A second ranger records \emph{Under 2 hours} / \emph{2 hours or more}. Which kind
is that?}{2.0cm}{Categorical --- two group labels; still grouping, just two bins} \\ \hline
\end{probgrid}
\\ \hline
% ── ROW 4: digits are not amounts — THE CRUX ─────────────────────────────────
\mainidea[Digits are not]{Amounts} &
A column of digits is \textbf{not automatically quantitative}: a permit number \emph{identifies}
a party and measures nothing. Counting a categorical variable also makes a number --- ``62\% hiked
with a dog'' is a \emph{summary} of \emph{Dog along}, not a value of it.

\vspace{4pt}
\begin{center}
\begin{tikzpicture}[scale=1.0]
  \foreach \x/\h/\a/\b/\c in {1.0/Permit \#/2041/6318/0377, 3.0/Parking spot/04/17/22,
                              5.0/Home ZIP/80211/80304/80027, 7.0/Water bottles/6/2/3,
                              9.0/Group size/4/2/5} {
    \fill[sky] (\x-0.95,0) rectangle (\x+0.95,2.4);
    \fill[navy] (\x-0.95,2.4) rectangle (\x+0.95,3.0);
    \draw[navy] (\x-0.95,0) rectangle (\x+0.95,3.0);
    \node[white, font=\bfseries\scriptsize] at (\x,2.7) {\h};
    \node[font=\normalsize] at (\x,1.95) {\a};
    \node[font=\normalsize] at (\x,1.2) {\b};
    \node[font=\normalsize] at (\x,0.45) {\c};}
  \node[font=\scriptsize\itshape] at (5.0,-0.35) {five columns, all digits --- not all amounts};
\end{tikzpicture}
\end{center}

\vspace{2pt}
\textbf{11.} Label each column \textbf{C} or \textbf{Q}:

\vspace{2pt}
\renewcommand{\arraystretch}{1.5}
\begin{tabularx}{\linewidth}{@{} Y Y Y Y Y @{}}
\textbf{Permit \#} & \textbf{Parking spot} & \textbf{Home ZIP} & \textbf{Water bottles} &
\textbf{Group size} \\
\hline
\ans{C} & \ans{C} & \ans{C} & \ans{Q} & \ans{Q} \\
\end{tabularx}
\\
& \notesprompt{Then say why.}
\begin{probgrid}{|Y|Y|}
\pcell{12}{Trail difficulty: \emph{Easy} / \emph{Moderate} / \emph{Hard}. Hard is more than Easy
--- so is it quantitative? How much more?}{2.2cm}{Categorical --- ordered labels, but no number of anything; ``how much more'' has no answer.} &
\pcell{13}{Parking spots 04 and 17. Average them, then hand me the party whose spot that is.
Which kind is \emph{Parking spot}?}{2.2cm}{$(4+17)/2 = 10.5$ --- nobody's spot. The digits name a spot, they measure nothing. C.} \\ \hline
\end{probgrid}
\\
& \begin{probgrid}*{|Y|Y|}
\pcell{14}{\emph{Dog along} is all Yes or No. Where did the 62 come from --- is it a value of
the variable?}{2.2cm}{From counting the Yes entries. No --- 62\% is a summary; the values are Yes and No, so the variable is categorical.} &
\pcell{15}{State the amount test in one sentence, without the word \emph{digits}.}{2.2cm}{Ask whether the value is an amount of something --- how much or how many --- or only a name for a group.} \\ \hline
\end{probgrid}
\\ \hline
% ── ROW 5: guided practice — WE DO, together, fresh context ──────────────────
\mainidea[Guided practice]{The Season Pass} &
A ski area records five things about each of its 3{,}400 pass holders. One holder's pass:

\vspace{4pt}
\begin{center}
\begin{tikzpicture}[scale=1.0]
  \fill[goldbg, rounded corners=6pt] (0,0) rectangle (8.6,3.4);
  \draw[goldacc, thick, rounded corners=6pt] (0,0) rectangle (8.6,3.4);
  \fill[navy, rounded corners=6pt] (0,2.6) rectangle (8.6,3.4);
  \fill[navy] (0,2.6) rectangle (8.6,2.9);
  \node[white, font=\bfseries] at (4.3,3.0) {SUMMIT RIDGE \quad SEASON PASS};
  \foreach \x/\y/\h/\v in {0.4/2.1/Pass number/5120, 4.5/2.1/Age (years)/34,
                           0.4/1.35/Pass tier/Gold, 4.5/1.35/Miles traveled/42,
                           0.4/0.6/Locker number/117} {
    \node[anchor=west, font=\scriptsize\color{slate}] at (\x,\y+0.25) {\h};
    \node[anchor=west, font=\bfseries\normalsize] at (\x,\y-0.1) {\v};}
  \draw[goldacc] (4.3,0.2) -- (4.3,2.5);
\end{tikzpicture}
\end{center}
\\
& \notesprompt{We work these together.}
\begin{probgrid}{|Y|Y|}
\pcell{16}{Who are the individuals? Label all five \textbf{C} or \textbf{Q}.}{2.4cm}{The 3{,}400 pass holders. Pass number C, Age Q, Pass tier C, Miles traveled Q, Locker number C.} &
\pcell{17}{\emph{Miles traveled} and \emph{Locker number} are both digits. Say the difference in
one sentence.}{2.4cm}{Miles measures how far (42 is more than 21); a locker number only names a locker (117 is not ``more locker'' than 58).} \\ \hline
\end{probgrid}
\\
& \begin{probgrid}*{|Y|Y|}
\pcell{18}{Age in years becomes \textbf{age group} (Under 18 / 18--34 / 35--54 / 55+). Which
kind now, and what can the resort no longer compute?}{2.4cm}{Categorical --- four group labels. The resort can no longer compute the mean age (or any exact age).} &
\pcell{19}{\textbf{``The average pass number is 5{,}120.''} Why does that tell you
nothing?}{2.4cm}{Pass numbers identify holders; averaging labels gives a number that is nobody's pass and measures nothing.} \\ \hline
\end{probgrid}
\\ \hline
\end{guidednotes}

% ── THE NOTES END AT GUIDED PRACTICE ─────────────────────────────────────────
% There is deliberately no "you do" here: ../homework/ is the individual practice,
% started in class in the last 8 minutes while the teacher circulates.

\end{document}
